\begin{abstract}
We introduce a two-branch benchmark for coherent memory in a sequential quantum process. A trusted verifier sends one half of each of $n$ independent EPR pairs through a device, requires every output to be committed before the next input is released, and sequesters those outputs. After the stream and the device's terminal-memory commitment, a hidden choice requests either (AUDIT) the parity of later computational-basis measurements on the references or (RETURN) joint recovery of all EPR pairs and visible reset of the terminal memory. For the explicit null class of arbitrary adaptive local instruments with unlimited classical memory and disposable within-slot ancillas, but no coherent state or entanglement persisting between slots, we prove the exact support
\[
\Cstream=max_{0\leq t\leq1}\left[
\frac{\lambda}{2}(1+t^n)+(1-\lambda)
\left(\frac{1+\sqrt{1-t^2}}2\right)^n\right].
\]
The reduction covers non-QND instruments and arbitrary transcript-conditioned joint recovery. At equal branch weight, the null is $3/4$ for every $n\geq2$, whereas one persistent coherent qubit attains $(\PA,\FR)=(1,1)$. A collective instrument retaining only a classical outcome has the strictly larger support
$\Ccoll=\tfrac12+\tfrac12\sqrt{\lambda^2+(1-\lambda)^2}$, isolating the causal cost of streaming. For $n\geq3$, the optimal classical strategy exhibits a first-order transition from recording nothing to an almost projective per-slot measurement. We supply a fixed-sample statistical certificate with explicit systematic allowances, conservative power planning, adversarial numerical checks, deterministic data and figures, and a hardware-facing protocol. The result is a trusted-interface temporal quantum-memory witness, not a device-independent claim, a proof of global erasure, or evidence for a preferred interpretation of quantum mechanics.
\end{abstract}

\noindent\textbf{Keywords:} quantum memory; multitime quantum processes; information--disturbance; entanglement fidelity; reversible computation; causal benchmarks.
